%=============================================================
\chapter{Performance Evaluation}
\label{ch:evaluation}
%=============================================================

本章評估 PaceANN 的兩個機制在伺服器模式下的效果。所有實驗沿用第~\ref{ch:method}~章的量測協定：
以 client--server 方式進行，原始查詢集只跑一次（不做 warmup、不重複查詢，除局部性實驗外也不套用 zipf 分布），
並在相同記憶體預算與相同召回率（iso-recall）下比較各系統的 QPS、平均延遲、尾延遲，以及每次查詢的平均 I/O 數。
比較對象包含 DiskANN baseline、Starling 與 PipeANN。並發度取 $C=8$（與 PipeANN 在本機的公平執行點對齊）與 $C=32$。

本版先固定實驗設計、欄位與比較方式；表中的實測數字待在實驗室機器上依相同配置重跑後填入，
標示為「待填」的欄位即為此類佔位。

%-------------------------------------------------------------
\section{iso-recall 下的系統比較}
\label{sec:eval-compare}
%-------------------------------------------------------------

第一組實驗在均勻（i.i.d.）查詢下，比較四個系統在相同 8\,GB 記憶體預算、相同召回率門檻下的效能。
各系統各自掃描搜尋長度 $L$，得到召回率對 QPS、尾延遲與 I/O 的 Pareto 曲線，再在共同的召回率點上比較。
PaceANN 在此採用 OPQ 壓縮（32\,bytes/vector）與 beam width 16 的配置，並同時啟用動態快取與早停。

\begin{table}[h]
\centering
\caption{iso-recall 下四系統的效能比較（i.i.d.，$C=8$，8\,GB 預算）}
\label{tab:main-compare}
\begin{tabular}{llccc}
\toprule
Recall@10 & 系統 & QPS & P99 (ms) & 平均 I/O \\
\midrule
$\approx 0.96$  & DiskANN  & 待填 & 待填 & 待填 \\
                & Starling & 待填 & 待填 & 待填 \\
                & PipeANN  & 待填 & 待填 & 待填 \\
                & PaceANN  & 待填 & 待填 & 待填 \\
\midrule
$\approx 0.97$  & DiskANN / Starling / PipeANN / PaceANN & 待填 & 待填 & 待填 \\
\midrule
$\approx 0.98$  & DiskANN / Starling / PipeANN / PaceANN & 待填 & 待填 & 待填 \\
\bottomrule
\end{tabular}
\end{table}

在高並發（$C=32$）下，我們另外比較各系統的 QPS 與尾延遲。PipeANN 在本機 16 執行緒下，
其 io\_uring 完成路徑需要 busy-poll、每個搜尋執行緒約佔一核，$C=32$ 時會超額訂閱而效能崩潰；
因此該點的公平比較改在執行緒較多的機器上進行，數字同樣待補（見 §\ref{sec:eval-discuss} 的公平性討論）。

%-------------------------------------------------------------
\section{兩機制的消融}
\label{sec:eval-ablation}
%-------------------------------------------------------------

為分別量化動態快取與早停的貢獻，我們在同一組 index 與記憶體預算下，
依序開關兩個機制，觀察 iso-recall 下的 QPS 與尾延遲變化。設定組合如第~\ref{ch:method}~章表~\ref{tab:ablation} 所列。

\begin{table}[h]
\centering
\caption{兩機制消融（i.i.d.，8\,GB 預算，固定召回率門檻）}
\label{tab:eval-ablation}
\begin{tabular}{lccc}
\toprule
設定 & QPS ($C=8$) & QPS ($C=32$) & 平均 I/O \\
\midrule
baseline（無快取無 ET） & 待填 & 待填 & 待填 \\
cache\_only            & 待填 & 待填 & 待填 \\
ET\_only               & 待填 & 待填 & 待填 \\
full                   & 待填 & 待填 & 待填 \\
\bottomrule
\end{tabular}
\end{table}

其中 \texttt{full} 對 \texttt{cache\_only} 的差距，是在相同 OPQ、beam width 與快取預算下\emph{只切換早停開關}所得，
因此可視為早停這一機制的乾淨貢獻。我們特別關注此貢獻是否隨並發度上升而增大：
若如此，代表早停省下的 I/O 在高並發時把頻寬讓回給佇列、進而放大吞吐量，這正是兩機制需要共同設計的直接證據。

%-------------------------------------------------------------
\section{工作負載局部性}
\label{sec:eval-locality}
%-------------------------------------------------------------

動態快取的價值取決於查詢是否具有局部性。我們以 Zipf 參數 $s$ 由 $0$（無重複、純均勻）
掃到 $1.5$（高度集中），觀察各系統的 QPS 與平均 I/O 隨局部性變化的趨勢。

\begin{table}[h]
\centering
\caption{局部性 sweep：QPS 與平均 I/O 隨 Zipf $s$ 的變化（固定召回率）}
\label{tab:eval-locality}
\begin{tabular}{lccc}
\toprule
Zipf $s$ & DiskANN QPS & PaceANN QPS & PaceANN 平均 I/O \\
\midrule
0.0（均勻） & 待填 & 待填 & 待填 \\
0.5        & 待填 & 待填 & 待填 \\
1.0        & 待填 & 待填 & 待填 \\
1.5        & 待填 & 待填 & 待填 \\
\bottomrule
\end{tabular}
\end{table}

預期的趨勢是：局部性越強，PaceANN 的動態快取命中越多、平均 I/O 塌陷、QPS 顯著上升；
靜態的 DiskANN 幾乎持平，而 PipeANN 因為沒有自適應快取，也無法從局部性受益。

%-------------------------------------------------------------
\section{高並發尾延遲}
\label{sec:eval-tail}
%-------------------------------------------------------------

尾延遲是伺服器情境的關鍵指標。我們固定召回率，把並發度 $C$ 由 $1$ 掃到機器的執行緒上限，
比較各系統的 P99 與 P999 延遲。

\begin{table}[h]
\centering
\caption{尾延遲隨並發度的變化（固定召回率）}
\label{tab:eval-tail}
\begin{tabular}{lcccc}
\toprule
並發度 $C$ & DiskANN P99 & Starling P99 & PipeANN P99 & PaceANN P99 \\
\midrule
1  & 待填 & 待填 & 待填 & 待填 \\
8  & 待填 & 待填 & 待填 & 待填 \\
16 & 待填 & 待填 & 待填 & 待填 \\
32 & 待填 & 待填 & 待填 & 待填 \\
\bottomrule
\end{tabular}
\end{table}

動態快取透過座標共存讓命中節點免去 re-rank 的回盤，是壓低 P99／P999 的主要來源；
早停則減少進入共享 NVMe 佇列的冗餘讀取，在高並發時緩解佇列排隊。

%-------------------------------------------------------------
\section{$\theta$ 參數敏感性}
\label{sec:eval-theta}
%-------------------------------------------------------------

早停的鬆緊由門檻 $\theta$ 控制。$\theta$ 越小越積極、省下的 hop 越多，但假陽性終止造成的召回率損失也越大。
我們在固定資料集與快取配置下掃描 $\theta$，觀察其對召回率與 QPS 的取捨，並確認最優區間的穩定性。

\begin{table}[h]
\centering
\caption{$\theta$ 對召回率與 QPS 的影響（$C=8$）}
\label{tab:eval-theta}
\begin{tabular}{lcc}
\toprule
$\theta$ & Recall@10 & QPS \\
\midrule
1.2 & 待填 & 待填 \\
1.3 & 待填 & 待填 \\
1.4 & 待填 & 待填 \\
1.6 & 待填 & 待填 \\
\bottomrule
\end{tabular}
\end{table}

$\theta$ 的最優點與資料集的 PQ 量化誤差分布有關，不同維度與 segment 數的 PQ 設定下略有差異，
因此本文將 $\theta$ 視為需依資料集校準的參數，而非固定常數。

%-------------------------------------------------------------
\section{討論與局限}
\label{sec:eval-discuss}
%-------------------------------------------------------------

綜合以上實驗，PaceANN 的兩個機制各自針對磁碟式搜尋在空間與時間上的浪費：
動態快取在有局部性的工作負載上把熱點留在 DRAM、消除 re-rank 回盤，是壓低尾延遲的主力；
早停則減少收斂後的冗餘 I/O，在高並發下把頻寬讓回佇列而提升吞吐量。兩者在伺服器情境下互補。

本研究仍有若干界限，於此如實列出。
第一，目前的主結果以 SpaceV 為主，尚需在 SIFT 與 DEEP 等資料集上複製，才能宣稱跨資料集的普遍性。
第二，i.i.d. 情境下的優勢部分來自 OPQ 壓縮與 beam width 16 這組共用配置；這是合法的設計選擇，但嚴格的公平比較應讓對手也採用相同壓縮，屆時早停仍是 PaceANN 獨有的貢獻。
第三，PipeANN 在本機的公平執行點為 $C=8$（8 個搜尋執行緒加 8 個 poller 恰好用滿 16 執行緒）；其在更高並發下的崩潰源自 busy-poll 的執行緒超額訂閱，而非索引品質，因此高並發的比較需在執行緒較多的機器上進行。
第四，早停的參數（$\theta$、patience、最小 hop 數）目前為靜態設定，未隨查詢動態調整。
第五，評估僅涵蓋 1 億級資料與單顆 SSD 上的 io\_uring，尚未驗證十億級規模、多 SSD 或 SPDK 等情境。
這些界限同時也是後續工作的方向。
